\documentclass[a4paper,10pt]{jsarticle}

% 英文化
\usepackage[english]{babel}
% 数式
\usepackage{amsmath}
\usepackage{amsfonts}
\usepackage{amsthm}
\usepackage{amssymb}
\usepackage{bm}
\usepackage{mathtools}
% 画像
\usepackage[dvipdfmx]{graphicx}
% 箇条書き
\usepackage{enumitem}

% 枠付き文章
\usepackage{ascmac}


\newcommand{\Af}{\footnotesize\it}
\newcommand{\Au}{\normalsize}
\newcommand{\Ti}{\LARGE}

\begin{document}

\begin{center}\Ti
  Set-Valued Fan-Takahashi Inequalities with Set-Relations Based on Scalarization Methods
\end{center}
\vskip 1em
\begin{center}\Au
  Ryota Iwamoto$^\ast$$^{a}$, Tamaki Tanaka$^{b}$
\end{center}
\begin{center}\Af
  $^{a}$Graduate School of Science and Technology,
  Niigata University, Niigata 950--2181, Japan.\\
  $^{b}$Faculty of Science, Niigata University, Niigata 950--2181, Japan.

  \vskip 1em
  $^\ast$presenting author.\\
  email: f24j003a@mail.cc.niigata-u.ac.jp
\end{center}
\date{}

In convex analysis and optimization theory,
Fan-Takahashi minimax inequality plays a key role to solve equilibrium problems and
other problems.
Let $X$ be a nonempty compact convex subset of a Hausdorff topological vector space and $f \colon X \times X \to \mathbb{R}$.
Fan-Takahashi minimax inequality is: if $f$ satisfies
the following conditions:
\begin{enumerate}[label=(\alph*)]
  \item for each fixed $y \in X$, $f(\cdot,y)$ is lower semicontinuous,
  \item for each fixed $x \in X$, $f(x,\cdot)$ is quasi concave,
  \item $f(x,x) \leq 0$ for all $x \in X$,
\end{enumerate}
then there exists $\bar{x} \in X$ such that $f(\bar{x},y) \leq 0$ for all $y \in X$.

About a quarter of a century ago, Georgiev and Tanaka [2] extended the minimax inequality to the form of set-valued maps.
After that, Kuwano, Tanaka, and Yamada [4] constructed the result of four types of set-valued minimax inequalities with set relations [3],
which are binary relations depending on a given convex cone. However, this result is limited to the case of specific scalarization functions. To obtain more practical results,
we need to replace them into more general scalarization functions. In addition,
Dechboon and Tanaka [1] proposed generalized
semi-continuity to inherit properties of cone continuity for set-valued maps.

The aim of this talk is to explain Fan-Takahashi minimax inequality and our results
in detail.

\begin{thebibliography}{99}
  \bibitem{} {\rm P.~Dechboon and T.~Tanaka},
  {\em Inheritance Properties on Cone Continuity for Set-Valued Maps via Scalarization},
  Minimax Theory and its Applications, {\bf 9(2)} (2024), 201--224.
  \bibitem{} {\rm P.G.~Georgiev and T.~Tanaka},
  {\em Vector-valued set-valued variants of Ky Fan's inequality},
  J. Nonlinear and Convex Anal., {\bf 1} (2000), 245--254.
  \bibitem{} {\rm D.~Kuroiwa, T.~Tanaka, and T.X.D.~Ha},
  {\em  On cone convexity of set-valued maps},
  Nonlinear Anal., {\bf 30} (1997) 1487--1496.
  \bibitem{} {\rm I.~Kuwano, T.~Tanaka, and S.~Yamada},
  {\em Unified scalarization for sets and set-valued Ky Fan minimax inequality},
  J. Nonlinear Convex Anal., {\bf 11} (2010), 513--525.
\end{thebibliography}

\end{document}